\input{\string~/Documents/School/"UC Irvine"/ta_template2.0.tex}
\rh{2B}{5.5}
\chead{\today}
\graphicspath{ {\string~/Desktop} }
\usetikzlibrary{arrows}
%============ANSWER KEY TOGGLE===========
\newcommand{\ans}[1]{ \noindent {\color{red} \textbf{Answer:} #1}}
%\renewcommand{\ans}[1]{\vspace{3in}}
%==========================================
\begin{document}
\begin{center}
\textbf{The Substitution Rule}
\end{center}

\begin{prop}[$u$-Substitution]
If $g'$ is a continuous on $[a,b]$ and $f$ is continuous on the range of $u = g(x)$, then:
\[
\int_a^b f(g(x))g'(x)dx =  \int_{g(a)}^{g(b)} f(u)du
\]
For indefinite integrals, we have:
\[
\int f(g(x))g'(x)dx = \int f(u)du
\]
\end{prop}

\begin{aun}[Symmetry]
For functions that are even or odd, we have little shortcuts that can help us evaluate integrals quickly:\\
\begin{center}
\begin{tabular}{|c|c|c|}
\hline
\textbf{Type} & \textbf{Definition} & \textbf{Integral Property} \\ \hline
$f(x)$ is Even & $f(-x) = f(x)$ & $\int_{-a}^a f(x) dx = 2\int_0^a f(x)dx$\\
$f(x)$ is Odd & $f(-x) = -f(x)$ & $\int_{-a}^a f(x) dx = 0$ \\ \hline
\end{tabular}
\end{center}
Use this as a trick to potentially bypass tricky problems!
\end{aun}


\begin{exx}
Compute the definite integral:
\[
\int_1^2 \frac{dx}{(3-5x)^2}
\]
\end{exx}

\ans{Let $u = 3 - 5x$}

\problembox{Compute the following definite integral:
\[
\int_e^{e^4} \frac{1}{x \sqrt{\ln(x)}}
\]}

\ans{Let $u = \ln(x)$}

\problembox{Compute the following indefinite integrals
\begin{enumerate}[label=$(\roman*)$]
\item $\int \sin^2 \theta \cos \theta d\theta$
\item $ \int \sec^2 2\theta d \theta$
\item $\int \frac{\cos (\ln(t))}{t} dt$
\item $\int \tan^2 \theta \sec^2 \theta d\theta $
\end{enumerate}
}

\ans{We'll list the $u$ choices (I'll complete one to make sure we know what we're doing)
\begin{enumerate}[label=$(\roman*)$]
\item $u = \sin(x), du = \cos(x)dx$ We can directly make a substitution:
\[
\int \sin^2 \theta \cos(\theta) d\theta = \int u^2 du = \frac{u^3}{3} + C
\]
This is an answer in terms of $u$, so to complete the problem, we must substitute back:
\[
\frac{u^3}{3} + C = \frac13 \sin^3 (x) + C
\]
completing our solution.
\item $u = 2\theta$
\item $u = \ln(t)$
\item $u = \tan(\theta)$
\end{enumerate}}

\problembox{[Challenge] Compute the following integral:
\[
\int_0^1 \frac{dx}{(1 + \sqrt{x})^4}
\]}

\ans{Final answer is $\frac16$. Here's a solution: Let $u = \sqrt{x} + 1$ and notice that $dx = 2(u-1)$ (rearrange the differentiated equation). This gives you the integral:
\[
\int_1^2 \frac{u-1}{u^4}du
\]
which can be split into two integrals, each of which are computable to us. 
}

\problembox{[Challenge] Compute the following indefinite integrals.
\begin{enumerate}[label=$\roman*.$]
\item $\displaystyle \int_{-\pi/2}^{\pi/2} \frac{x^2 \sin(x)}{x^6 + 1} dx $
\item $ \displaystyle \int \sqrt{x-1}(x+3) dx $
\item $\displaystyle \int \sqrt{x} \sin(\sqrt{x^3} +1)dx $
\end{enumerate}
}

\ans{ The first is 0 since it is an odd function. Second, substitute $u = x -1$; then $x+ 3 = u + 4$ and $du = dx$. Finally, the last one is $u = \sqrt{x^3} + 1$ and $du = \frac32 \sqrt{x}$
}

\problembox{[Challenge] Suppose $f$ is a function with $\displaystyle \int_0^8 f(x)dx = 20$. What is the value of $\displaystyle \int_0^2 f(4x)dx$?}

\ans{Using $u = 4x$, we see that:
\[
\int_0^2 f(4x) dx = \frac14 \int_0^8 f(u) du = \frac14 \cdot 20 = 5
\]}

\end{document}